% !TEX root = ../高等概率论.tex
% =====================================================================
%  弱大数定律：L² 弱大数定律，三角阵列，截断法，Khinchin 弱大数定律，St. Petersburg 悖论
% =====================================================================
\lecture{不相关随机变量的 \texorpdfstring{$L^2$}{L²} 弱大数定律，Bernstein 多项式与高维立方体，三角阵列与集券问题、随机置换的圈、占位问题，截断法与三角阵列的弱大数定律，Khinchin 弱大数定律，Cauchy 分布与 St.\ Petersburg 悖论}
\label{lec:weak-laws}

把一枚公平的硬币抛 $n$ 次，正面的次数 $S_n$ 满足 $\E S_n=\frac n2$，$\Var(S_n)=\frac n4$。由 Chebyshev 不等式，
\[
\P\Bigl(\Bigl\lvert\frac{S_n}n-\frac12\Bigr\rvert\geqslant\varepsilon\Bigr)\leqslant\frac{\Var(S_n)}{n^2\varepsilon^2}=\frac1{4n\varepsilon^2}\to0 .
\]
这就是第~\ref{lec:independence} 讲的 Bernoulli 弱大数定律（命题~\ref{prop:independence:bernoulli}）。这个证明只用到了两件事：$\Var(S_n)$ 只有 $n$ 的量级，以及 Chebyshev 不等式。这一讲沿着这条思路走下去，回答三个问题。
\begin{enumerate}
\item 硬币换成一般的随机变量 $X_1,X_2,\dots$，需要什么条件？第一节的回答是：只要两两不相关、方差有界，证明一字不改。
\item 若随机变量不是同一个序列，而是每个 $n$ 各有一组（比如集齐 $n$ 种卡片所需的次数），怎么办？第二节的回答是：只要方差相对于“尺度”$b_n^2$ 是小量。
\item 若方差无穷，甚至期望无穷，还有没有大数定律？第三节用截断法回答：i.i.d.\ 时，只要 $\E\abs{X_1}<\infty$，就有 $S_n/n\toP \E X_1$；甚至 $\E\abs{X_1}=\infty$ 时，有时也能找到合适的“中心”$a_n$ 使 $S_n/n-a_n\toP0$，St.\ Petersburg 悖论就是这样解决的。
\end{enumerate}
这一讲的结论都是\term{弱大数定律}（weak law of large numbers），结论是依概率收敛。几乎必然收敛的\term{强大数定律}在后面的一讲中证明。

\topic{本讲用到的前文结果}

\begin{enumerate}
\item 依概率收敛（定义~\ref{def:convergence:probability}）：$Y_n\toP Y$ 指对一切 $\varepsilon>0$，$\P(\abs{Y_n-Y}>\varepsilon)\to0$。$L^p$ 收敛蕴含依概率收敛（定理~\ref{thm:convergence:Lp-implies-p}）：由 Chebyshev 不等式，$\P(\abs{Y_n-Y}\geqslant\varepsilon)\leqslant\varepsilon^{-p}\E\abs{Y_n-Y}^p$。依概率收敛对连续函数与四则运算封闭（推论~\ref{cor:convergence:continuous-mapping}）。
\item 方差与协方差（定义~\ref{def:expectation:variance}、命题~\ref{prop:expectation:variance}）：$\Var(aX+b)=a^2\Var(X)$；$\Var(\sum X_i)=\sum\Var(X_i)+2\sum_{i<j}\Cov(X_i,X_j)$；$\Var(X)\leqslant \E X^2$。
\item Chebyshev 不等式（推论~\ref{cor:expectation:chebyshev}）：$\P(\abs{X-\E X}\geqslant a)\leqslant\Var(X)/a^2$。
\item 独立随机变量两两不相关，和的方差等于方差的和（推论~\ref{cor:independence:variance}）；独立随机变量的函数仍独立（推论~\ref{cor:independence:functions}）；离散随机变量的独立性只需验证联合概率质量函数是乘积（推论~\ref{cor:independence:discrete}）；独立随机变量序列存在（定理~\ref{thm:independence:existence}）。
\item 尾概率公式（定理~\ref{thm:product:tail-formula}）：$Y\geqslant0$，$p>0$ 时，$\E Y^p=\int_0^\infty py^{p-1}\P(Y>y)\dd y$。
\item 课堂提问的例题（第~\pageref{we:expectation:tP} 页）：$\E\abs X<\infty$ 蕴含 $x\P(\abs X>x)\to0$，而且 $\E(\abs X;\abs X>x)\to0$。
\item 控制收敛定理与有界收敛定理（定理~\ref{thm:integration:dct}、推论~\ref{cor:integration:bct}）。
\end{enumerate}

\topic{\texorpdfstring{$L^2$}{L²} 弱大数定律}

若一族随机变量 $\{X_i\}_{i\in I}$ 满足 $\E X_i^2<\infty$，并且 $\E(X_iX_j)=\E X_i\,\E X_j$（即 $\Cov(X_i,X_j)=0$）对一切 $i\neq j$ 成立，称它们\term{两两不相关}（uncorrelated）。独立且平方可积的随机变量两两不相关（推论~\ref{cor:independence:variance}），反之不然（\egref{ex:independence:uncorrelated}）。

\begin{theorem}\label{thm:weak-laws:variance-sum}
设 $X_1,\dots,X_n$ 平方可积且两两不相关。则 $\Var(X_1+\cdots+X_n)=\Var(X_1)+\cdots+\Var(X_n)$。
\end{theorem}
\begin{proof}
这是命题~\ref{prop:expectation:variance}(4) 的特例：右边多出的那一项 $2\sum_{i<j}\Cov(X_i,X_j)$ 为零。我们再把计算写一遍，因为它是本讲一切估计的出发点。记 $\mu_i=\E X_i$，$S_n=\sum_iX_i$，则 $\E S_n=\sum_i\mu_i$，
\[
\Var(S_n)=\E\Bigl(\sum_{i=1}^n(X_i-\mu_i)\Bigr)^2=\sum_{i=1}^n\E(X_i-\mu_i)^2+2\sum_{1\leqslant j<i\leqslant n}\E\bigl[(X_i-\mu_i)(X_j-\mu_j)\bigr],
\]
而 $\E[(X_i-\mu_i)(X_j-\mu_j)]=\E X_iX_j-\mu_i\E X_j-\mu_j\E X_i+\mu_i\mu_j=\E X_iX_j-\mu_i\mu_j=0$。
\end{proof}

\begin{theorem}[$L^2$ 弱大数定律]\label{thm:weak-laws:L2}
设 $X_1,X_2,\dots$ 两两不相关，$\E X_i=\mu$，$\Var(X_i)\leqslant C<\infty$。令 $S_n=X_1+\cdots+X_n$。则 $n\to\infty$ 时 $S_n/n\to\mu$ 在 $L^2$ 中成立，从而 $S_n/n\toP\mu$。
\end{theorem}
\begin{proof}
$\E(S_n/n)=\mu$，由 $\Var(cY)=c^2\Var(Y)$ 与定理~\ref{thm:weak-laws:variance-sum}，
\[
\E\Bigl(\frac{S_n}n-\mu\Bigr)^2=\Var\Bigl(\frac{S_n}n\Bigr)=\frac1{n^2}\bigl(\Var(X_1)+\cdots+\Var(X_n)\bigr)\leqslant\frac{Cn}{n^2}=\frac Cn\to0 .
\]
再由定理~\ref{thm:convergence:Lp-implies-p}（$p=2$），$S_n/n\toP\mu$。具体地，$\P(\abs{S_n/n-\mu}>\varepsilon)\leqslant\frac C{n\varepsilon^2}$。
\end{proof}

\begin{remark}
这个定理最重要的特例是：$X_1,X_2,\dots$ \term{独立同分布}（independent and identically distributed，简写为 i.i.d.），$\E X_1^2<\infty$，则 $S_n/n\toP \E X_1$。定理~\ref{thm:weak-laws:khinchin} 将说明 $\E\abs{X_1}<\infty$ 就够了。但 $L^2$ 的证明有它的好处：它给出了收敛的速度 $C/(n\varepsilon^2)$，并且只用到“不相关”，远比独立弱。证明的要点是一个数量级的比较：$S_n-\E S_n$ 的“典型大小”是标准差 $\sqrt{\Var(S_n)}\leqslant\sqrt{Cn}$，比 $n$ 小得多，所以除以 $n$ 以后趋于零。以后的中心极限定理会说，$\sqrt n$ 恰好是正确的量级。
\end{remark}

我们先看两个看起来与随机性无关的应用。

\begin{workedexample}[多项式逼近]\label{we:weak-laws:bernstein}
设 $f$ 是 $[0,1]$ 上的连续函数，
\[
f_n(x)=\sum_{m=0}^n\binom nmx^m(1-x)^{n-m}f\Bigl(\frac mn\Bigr)
\]
是 $f$ 的 $n$ 次 \term{Bernstein 多项式}。证明 $\sup_{x\in[0,1]}\abs{f_n(x)-f(x)}\to0$。
\end{workedexample}
\begin{solution}
这正是第~\ref{lec:expectation} 讲作业 E 的 Weierstrass 定理，那里用的是同一个想法，这里换成大数定律的语言再说一遍。

\emph{第一步：概率解释。}固定 $p\in[0,1]$，取独立的 $X_1,\dots,X_n$，$\P(X_i=1)=p$，$\P(X_i=0)=1-p$（由定理~\ref{thm:independence:existence} 存在），$S_n=X_1+\cdots+X_n$。$S_n$ 服从二项分布，$\P(S_n=m)=\binom nmp^m(1-p)^{n-m}$，所以由变量替换公式（推论~\ref{cor:expectation:discrete}），$\E f(S_n/n)=f_n(p)$。直观上，$S_n/n\approx p$，所以 $f_n(p)=\E f(S_n/n)\approx f(p)$。

\emph{第二步：定量。}$\E X_i=p$，$\Var(X_i)=p(1-p)\leqslant\frac14$，由定理~\ref{thm:weak-laws:L2} 的证明，
\[
\P(\abs{S_n/n-p}>\delta)\leqslant\frac{p(1-p)}{n\delta^2}\leqslant\frac1{4n\delta^2},
\]
这个界与 $p$ 无关。

\emph{第三步：一致连续性。}令 $M=\sup_{[0,1]}\abs f<\infty$。给定 $\varepsilon>0$，由 $f$ 在 $[0,1]$ 上一致连续，取 $\delta>0$ 使 $\abs{x-y}\leqslant\delta$ 时 $\abs{f(x)-f(y)}<\varepsilon$。在事件 $\{\abs{S_n/n-p}\leqslant\delta\}$ 上，$\abs{f(S_n/n)-f(p)}<\varepsilon$；在其余地方，它不超过 $2M$。所以
\[
\abs{f_n(p)-f(p)}\leqslant \E\abs{f(S_n/n)-f(p)}\leqslant\varepsilon+2M\P(\abs{S_n/n-p}>\delta)\leqslant\varepsilon+\frac{2M}{4n\delta^2}.
\]
右边与 $p$ 无关，所以 $\limsup_n\sup_p\abs{f_n(p)-f(p)}\leqslant\varepsilon$，$\varepsilon$ 任意。

注意这里不能只用定理的结论“$S_n/n\toP p$”，而要用证明中的定量估计：一致收敛要求估计与 $p$ 无关。
\end{solution}

\begin{workedexample}[高维立方体几乎全在一个球壳上]\label{we:weak-laws:cube}
设 $X_1,X_2,\dots$ 独立，都服从 $(-1,1)$ 上的均匀分布。证明 $(X_1^2+\cdots+X_n^2)/n\toP\frac13$。由此说明：对任意 $\varepsilon>0$，记
\[
A_{n,\varepsilon}=\Bigl\{x\in\R^n:\ (1-\varepsilon)\sqrt{n/3}<\abs x<(1+\varepsilon)\sqrt{n/3}\Bigr\},
\]
则 $\lambda_n(A_{n,\varepsilon}\cap(-1,1)^n)/2^n\to1$。
\end{workedexample}
\begin{solution}
\emph{大数定律。}$Y_i=X_i^2$ 是独立随机变量的函数，仍独立（推论~\ref{cor:independence:functions}），同分布。$\E Y_i=\frac12\int_{-1}^1x^2\dd x=\frac13$，$\Var(Y_i)\leqslant \E Y_i^2=\frac12\int_{-1}^1x^4\dd x=\frac15\leqslant1$。由定理~\ref{thm:weak-laws:L2}，$(Y_1+\cdots+Y_n)/n\toP\frac13$。

\emph{几何解释。}$(X_1,\dots,X_n)$ 的分布是 $(-1,1)^n$ 上的均匀分布：由独立性与定理~\ref{thm:independence:product-measure}，它的分布是 $n$ 个一维均匀分布的乘积测度，即 $2^{-n}\lambda_n$ 限制在 $(-1,1)^n$ 上。所以
\[
\begin{aligned}
\frac{\lambda_n(A_{n,\varepsilon}\cap(-1,1)^n)}{2^n}&=\P\bigl((X_1,\dots,X_n)\in A_{n,\varepsilon}\bigr)\\
&=\P\Bigl((1-\varepsilon)^2\frac13<\frac{X_1^2+\cdots+X_n^2}n<(1+\varepsilon)^2\frac13\Bigr).
\end{aligned}
\]
令 $\eta=\frac13\min\{1-(1-\varepsilon)^2,(1+\varepsilon)^2-1\}>0$，则右边的事件包含 $\{\abs{\frac1n\sum Y_i-\frac13}<\eta\}$，其概率趋于 $1$。

这个结论很违反低维的直观：$n$ 很大时，立方体 $(-1,1)^n$ 的体积几乎全部集中在半径约为 $\sqrt{n/3}$ 的一个薄球壳附近。立方体的“角”到中心的距离是 $\sqrt n$，面的中心到中心的距离是 $1$，但这两类点都只占体积中可以忽略的一部分。这种“高维中测度集中”的现象，在统计学和机器学习中处处可见。
\end{solution}

\topic{三角阵列}

很多经典的极限定理研究的是\term{三角阵列}（triangular array）：对每个 $n$ 有一组随机变量 $X_{n,1},\dots,X_{n,n}$，考察行和 $S_n=X_{n,1}+\cdots+X_{n,n}$。不同的行可以完全不同。下面这个简单的定理连这个结构都不需要，$S_n$ 可以是任意一列随机变量。

\begin{theorem}\label{thm:weak-laws:triangular-L2}
设 $\mu_n=\E S_n$，$\sigma_n^2=\Var(S_n)$，$b_n>0$。若 $\sigma_n^2/b_n^2\to0$，则
\[
\frac{S_n-\mu_n}{b_n}\toP0 .
\]
\end{theorem}
\begin{proof}
$\E\bigl((S_n-\mu_n)/b_n\bigr)^2=b_n^{-2}\Var(S_n)\to0$，由定理~\ref{thm:convergence:Lp-implies-p}。
\end{proof}

定理的用法是：选一个尽量小、但仍使 $\sigma_n/b_n\to0$ 的尺度 $b_n$，结论就说 $S_n$ 在 $\mu_n$ 附近、误差为 $o(b_n)$。下面三个例子中要用到调和级数的估计：比较 $\frac1m$ 与 $\int_{m-1}^m\frac{\dd x}x$、$\int_m^{m+1}\frac{\dd x}x$，得
\begin{equation}\label{eq:weak-laws:harmonic}
\log n\leqslant\sum_{m=1}^n\frac1m\leqslant1+\log n .
\end{equation}
（左边：$\sum_{m=1}^n\frac1m\geqslant\int_1^{n+1}\frac{\dd x}x\geqslant\log n$；右边：$\sum_{m=2}^n\frac1m\leqslant\int_1^n\frac{\dd x}x=\log n$。）还要用到几何分布的方差：若 $\P(N=k)=p(1-p)^{k-1}$，$k\geqslant1$，则 $\E N=\frac1p$，$\Var(N)=\frac{1-p}{p^2}\leqslant\frac1{p^2}$。（由本讲习题~\exref{ex:weak-laws:A4}，$\E N^2=\sum_{k\geqslant1}(2k-1)\P(N\geqslant k)=\sum_{k\geqslant1}(2k-1)(1-p)^{k-1}=\frac2{p^2}-\frac1p$，减去 $(\E N)^2=\frac1{p^2}$。）

\begin{workedexample}[集券问题]\label{we:weak-laws:coupon}
设 $X_1,X_2,\dots$ 独立，都服从 $\{1,\dots,n\}$ 上的均匀分布。想象每买一包零食得到一张卡片，共有 $n$ 种，每次得到哪一种是等可能的、互相独立的。令 $\tau_k^n=\inf\{m:\ \abs{\{X_1,\dots,X_m\}}=k\}$ 为第一次集到 $k$ 种不同卡片的时刻，$T_n=\tau_n^n$ 为集齐全套的时刻。证明
\[
\frac{T_n-n\sum_{m=1}^n\frac1m}{n\log n}\toP0,\qquad\text{从而}\qquad\frac{T_n}{n\log n}\toP1 .
\]
\end{workedexample}
\begin{solution}
\emph{第一步：分阶段。}令 $\tau_0^n=0$，$X_{n,k}=\tau_k^n-\tau_{k-1}^n$（$1\leqslant k\leqslant n$）为“已有 $k-1$ 种之后，等到第 $k$ 种新卡片”所需的次数，$T_n=\sum_kX_{n,k}$。直观上，在第 $k$ 阶段每次抽到新卡片的概率是 $p_k=1-\frac{k-1}n$，所以 $X_{n,k}$ 服从参数 $p_k$ 的几何分布，并且各阶段互不影响。我们把这个直观证明出来。

\emph{第二步：联合分布。}固定正整数 $j_1,\dots,j_n$（必然 $j_1=1$），令 $L=j_1+\cdots+j_n$。事件 $\{X_{n,1}=j_1,\dots,X_{n,n}=j_n\}$ 只依赖于 $(X_1,\dots,X_L)$：它要求在第 $k$ 阶段的前 $j_k-1$ 次都抽到已有的 $k-1$ 种之一，第 $j_k$ 次抽到新的一种。满足这个要求的 $L$ 元组 $(x_1,\dots,x_L)\in\{1,\dots,n\}^L$ 的个数，按时间顺序逐个计数，是
\[
\prod_{k=1}^n(k-1)^{j_k-1}(n-k+1)
\]
（第 $k$ 阶段中每次“旧卡片”有 $k-1$ 种选法，“新卡片”有 $n-k+1$ 种选法）。由独立性与均匀性，每个 $L$ 元组的概率都是 $n^{-L}$。所以
\[
\P(X_{n,1}=j_1,\dots,X_{n,n}=j_n)=\prod_{k=1}^n\Bigl(\frac{k-1}n\Bigr)^{j_k-1}\frac{n-k+1}n=\prod_{k=1}^n(1-p_k)^{j_k-1}p_k .
\]
由推论~\ref{cor:independence:discrete}，$X_{n,1},\dots,X_{n,n}$ 独立，$X_{n,k}$ 服从参数为 $p_k$ 的几何分布。

\emph{第三步：期望与方差。}由线性性与 \eqref{eq:weak-laws:harmonic}，
\[
\E T_n=\sum_{k=1}^n\frac1{p_k}=\sum_{k=1}^n\frac n{n-k+1}=n\sum_{m=1}^n\frac1m\sim n\log n .
\]
由独立性（定理~\ref{thm:weak-laws:variance-sum}）与几何分布的方差，
\[
\Var(T_n)\leqslant\sum_{k=1}^n\frac1{p_k^2}=n^2\sum_{m=1}^n\frac1{m^2}\leqslant n^2\sum_{m=1}^\infty\frac1{m^2}=\frac{\pi^2}6n^2 .
\]

\emph{第四步：用定理~\ref{thm:weak-laws:triangular-L2}。}取 $b_n=n\log n$，$\Var(T_n)/b_n^2\leqslant\frac{\pi^2}{6(\log n)^2}\to0$，于是 $(T_n-\E T_n)/(n\log n)\toP0$。又 $\E T_n/(n\log n)\to1$（由 \eqref{eq:weak-laws:harmonic}），所以 $T_n/(n\log n)\toP1$。

例如 $n=365$：要“见到所有的生日”，大约需要 $365\log365\approx2153$ 个人，是生日种数的 $5.9$ 倍。$T_n$ 的波动 $\sqrt{\Var(T_n)}$ 是 $n$ 的量级，比 $n\log n$ 小一个对数因子，所以选 $b_n=n\log n$ 正好。
\end{solution}

\begin{workedexample}[随机置换的圈数]\label{we:weak-laws:cycles}
设 $\Omega_n$ 为 $\{1,\dots,n\}$ 的全体 $n!$ 个置换，赋予等可能的概率。把置换 $\pi$ 分解成圈：从 $1$ 出发看 $1,\pi(1),\pi(\pi(1)),\dots$，直到回到 $1$，这是第一个圈；再从不在第一个圈中的最小的数出发，如此下去。比如
\[
\begin{array}{c|ccccccccc}i&1&2&3&4&5&6&7&8&9\\\hline\pi(i)&3&9&6&8&2&1&5&4&7\end{array}
\]
的圈分解是 $(1\,3\,6)(2\,9\,7\,5)(4\,8)$。按圈分解中出现的顺序写下这 $n$ 个数，令 $X_{n,k}=1$ 若第 $k$ 个数后面紧跟一个右括号，否则为 $0$；$S_n=X_{n,1}+\cdots+X_{n,n}$ 是圈的个数。（例子中 $X_{9,3}=X_{9,7}=X_{9,9}=1$。）证明
\begin{enumerate}
\item $X_{n,1},\dots,X_{n,n}$ 独立，$\P(X_{n,k}=1)=\frac1{n-k+1}$；
\item 对任意 $\varepsilon>0$，$\dfrac{S_n-\sum_{m=1}^n\frac1m}{(\log n)^{1/2+\varepsilon}}\toP0$，特别地 $S_n/\log n\toP1$。
\end{enumerate}
\end{workedexample}
\begin{solution}
\emph{第一步：边走边造一个置换。}我们用下面的办法逐步决定 $\pi$：令 $i_1=1$，在 $\{1,\dots,n\}$ 中选一个数 $j_1$，令 $\pi(i_1)=j_1$。一般地，设 $i_1,j_1,\dots,i_{k-1},j_{k-1}$ 已经选定，$\pi(i_l)=j_l$。若 $j_{k-1}\in\{i_1,\dots,i_{k-1}\}$，说明刚刚完成了一个圈，令 $i_k$ 为不在 $\{i_1,\dots,i_{k-1}\}$ 中的最小的数；否则令 $i_k=j_{k-1}$（沿着圈继续走）。然后在 $\{1,\dots,n\}\setminus\{j_1,\dots,j_{k-1}\}$ 这 $n-k+1$ 个数中选一个 $j_k$，令 $\pi(i_k)=j_k$。

每个选择序列 $(j_1,\dots,j_n)$ 给出一个置换，并且 $i_1,\dots,i_n$ 恰好是圈分解中依次写下的那些数。反过来，每个置换 $\pi$ 由它的圈分解唯一地对应一个选择序列（$j_k=\pi(i_k)$）。所以这是选择序列与置换之间的一一对应。选择序列共有 $n\cdot(n-1)\cdots1=n!$ 个，在等可能的置换下，每个选择序列的概率是 $\frac1{n!}=\prod_{k=1}^n\frac1{n-k+1}$。

\emph{第二步：右括号。}第 $k$ 个数 $i_k$ 后面是右括号，当且仅当 $\pi(i_k)=j_k$ 等于当前这个圈的起点 $c_k$。$c_k$ 由 $j_1,\dots,j_{k-1}$ 决定，并且 $c_k$ 还没有被选为任何 $j_l$（$l<k$），因为当前的圈还没有闭合。所以在第 $k$ 步的 $n-k+1$ 个可选的数中，恰好有一个（即 $c_k$）使 $X_{n,k}=1$，另外 $n-k$ 个使 $X_{n,k}=0$。

\emph{第三步：独立性。}给定 $e_1,\dots,e_n\in\{0,1\}$，满足 $X_{n,k}=e_k$（$k=1,\dots,n$）的选择序列的个数，按 $k$ 的顺序计数，是 $\prod_k c(e_k)$，其中 $c(1)=1$，$c(0)=n-k$。所以
\[
\begin{aligned}
&\P(X_{n,1}=e_1,\dots,X_{n,n}=e_n)=\prod_{k=1}^n\frac{c(e_k)}{n-k+1}=\prod_{k=1}^nq_k(e_k),\\
&\text{其中}\ q_k(1)=\frac1{n-k+1},\quad q_k(0)=\frac{n-k}{n-k+1}.
\end{aligned}
\]
由推论~\ref{cor:independence:discrete}，(1) 成立。（$k=n$ 时只剩一个数，必然闭合，$\P(X_{n,n}=1)=1$，与公式一致。）

\emph{第四步：大数定律。}由 (1)，$\E S_n=\sum_{k=1}^n\frac1{n-k+1}=\sum_{m=1}^n\frac1m$。$X_{n,k}^2=X_{n,k}$，所以
\[
\Var(S_n)=\sum_k\Var(X_{n,k})\leqslant\sum_k\E X_{n,k}^2=\sum_k\E X_{n,k}=\E S_n\leqslant1+\log n .
\]
取 $b_n=(\log n)^{1/2+\varepsilon}$，$\Var(S_n)/b_n^2\leqslant(1+\log n)/(\log n)^{1+2\varepsilon}\to0$，由定理~\ref{thm:weak-laws:triangular-L2} 得 (2)。取 $\varepsilon=\frac12$，再由 \eqref{eq:weak-laws:harmonic}，$S_n/\log n\toP1$。

所以一个随机置换大约有 $\log n$ 个圈。结论 (2) 比 $S_n/\log n\toP1$ 精确得多：$S_n$ 与 $\log n$ 的偏差不超过 $(\log n)^{1/2+\varepsilon}$。以后的中心极限定理会说明，$\varepsilon=0$ 时结论不成立，偏差恰好是 $\sqrt{\log n}$ 的量级。
\end{solution}

\begin{workedexample}[占位问题]\label{we:weak-laws:occupancy}
把 $r$ 个球随机地放进 $n$ 个盒子，即 $n^r$ 种放法等可能（每个球独立地、等可能地选一个盒子）。令 $N_n$ 为空盒子的个数。证明：若 $r=r(n)$ 满足 $r/n\to c\in(0,\infty)$，则 $N_n/n\toP e^{-c}$。
\end{workedexample}
\begin{solution}
令 $A_i$ 为“第 $i$ 个盒子是空的”，$N_n=\sum_{i=1}^n\one_{A_i}$。$A_i$ 要求 $r$ 个球都不放进第 $i$ 个盒子；$i\neq j$ 时 $A_i\cap A_j$ 要求 $r$ 个球都不放进这两个盒子。所以
\[
\P(A_i)=\Bigl(1-\frac1n\Bigr)^r,\qquad \P(A_i\cap A_j)=\Bigl(1-\frac2n\Bigr)^r\quad(i\neq j).
\]

\emph{期望。}$\E N_n=n(1-\frac1n)^r$。由 $\log(1-\frac1n)=-\frac1n+O(\frac1{n^2})$，$r\log(1-\frac1n)=-\frac rn+O(\frac r{n^2})\to-c$，所以 $(1-\frac1n)^r\to e^{-c}$，$\E N_n/n\to e^{-c}$。同理 $(1-\frac2n)^r\to e^{-2c}$。

\emph{方差。}$\E N_n^2=\E\bigl(\sum_i\one_{A_i}\bigr)^2=\sum_{i,j}\P(A_i\cap A_j)$，于是
\[
\begin{aligned}
\Var(N_n)&=\sum_{i,j}\bigl(\P(A_i\cap A_j)-\P(A_i)\P(A_j)\bigr)\\
&=n(n-1)\Bigl[\Bigl(1-\frac2n\Bigr)^r-\Bigl(1-\frac1n\Bigr)^{2r}\Bigr]+n\Bigl[\Bigl(1-\frac1n\Bigr)^r-\Bigl(1-\frac1n\Bigr)^{2r}\Bigr],
\end{aligned}
\]
前一项来自 $i\neq j$，后一项来自 $i=j$。除以 $n^2$：方括号中第一项趋于 $e^{-2c}-e^{-2c}=0$，第二项除以 $n^2$ 后不超过 $\frac1n$。所以 $\Var(N_n/n)\to0$。

\emph{结论。}取 $b_n=n$，由定理~\ref{thm:weak-laws:triangular-L2}，$(N_n-\E N_n)/n\toP0$，再结合 $\E N_n/n\to e^{-c}$ 即得。注意 $\one_{A_i}$ 之间并不独立（甚至是负相关的），这里只用到了方差的计算。比如 $r=n$ 时，大约有 $e^{-1}\approx37\%$ 的盒子是空的。
\end{solution}

\topic{截断法}

到目前为止的证明都靠方差。可是 i.i.d.\ 的随机变量只要 $\E\abs{X_1}<\infty$，大数定律就应该成立，而方差可能是无穷。处理的办法是\term{截断}（truncation）：把 $X$ 在水平 $M$ 处截断，即考虑
\[
\bar X=X\one_{\{\abs X\leqslant M\}}=\begin{cases}X,&\abs X\leqslant M,\\0,&\abs X>M.\end{cases}
\]
$\bar X$ 有界，方差有限，可以用 Chebyshev 不等式；而 $X$ 与 $\bar X$ 只在 $\{\abs X>M\}$ 上不同，这个事件的概率很小。截断的水平要选得恰当：太低，$X\neq\bar X$ 的概率太大；太高，$\bar X$ 的方差太大。下面的定理把这两方面的要求写成条件 (i)(ii)，它的证明是“假定了我们需要的东西”，所以很简单；难处在于对具体的例子验证这两个条件。

\begin{theorem}[三角阵列的弱大数定律]\label{thm:weak-laws:truncated-triangular}
对每个 $n$，设 $X_{n,1},\dots,X_{n,n}$ 独立。设 $b_n>0$，$b_n\to\infty$，令 $\bar X_{n,k}=X_{n,k}\one_{\{\abs{X_{n,k}}\leqslant b_n\}}$。若 $n\to\infty$ 时
\begin{enumerate}
\item[(i)] $\sum_{k=1}^n\P(\abs{X_{n,k}}>b_n)\to0$，
\item[(ii)] $b_n^{-2}\sum_{k=1}^n\E\bar X_{n,k}^2\to0$，
\end{enumerate}
令 $S_n=X_{n,1}+\cdots+X_{n,n}$，$a_n=\sum_{k=1}^n\E\bar X_{n,k}$，则 $(S_n-a_n)/b_n\toP0$。
\end{theorem}
\begin{proof}
令 $\bar S_n=\bar X_{n,1}+\cdots+\bar X_{n,n}$。若 $S_n\neq\bar S_n$，则某个 $X_{n,k}\neq\bar X_{n,k}$。所以对 $\varepsilon>0$，
\[
\P\Bigl(\Bigl\lvert\frac{S_n-a_n}{b_n}\Bigr\rvert>\varepsilon\Bigr)\leqslant \P(S_n\neq\bar S_n)+\P\Bigl(\Bigl\lvert\frac{\bar S_n-a_n}{b_n}\Bigr\rvert>\varepsilon\Bigr).
\]
\emph{第一项。}由次可加性与 (i)，
\[
\P(S_n\neq\bar S_n)\leqslant \P\Bigl(\bigcup_{k=1}^n\{\bar X_{n,k}\neq X_{n,k}\}\Bigr)\leqslant\sum_{k=1}^n\P(\abs{X_{n,k}}>b_n)\to0 .
\]
\emph{第二项。}$a_n=\E\bar S_n$。$\bar X_{n,k}$ 是独立随机变量的函数，仍独立（推论~\ref{cor:independence:functions}），有界，所以两两不相关。由 Chebyshev 不等式、定理~\ref{thm:weak-laws:variance-sum} 与 $\Var(Y)\leqslant \E Y^2$，
\[
\P\Bigl(\Bigl\lvert\frac{\bar S_n-a_n}{b_n}\Bigr\rvert>\varepsilon\Bigr)\leqslant\frac{\Var(\bar S_n)}{\varepsilon^2b_n^2}=\frac1{\varepsilon^2b_n^2}\sum_{k=1}^n\Var(\bar X_{n,k})\leqslant\frac1{\varepsilon^2b_n^2}\sum_{k=1}^n\E\bar X_{n,k}^2\to0
\]
由 (ii)。
\end{proof}

现在把它用到一个序列上。

\begin{theorem}[弱大数定律]\label{thm:weak-laws:wlln}
设 $X_1,X_2,\dots$ 是 i.i.d.，并且
\begin{equation}\label{eq:weak-laws:xP}
x\P(\abs{X_1}>x)\to0\qquad(x\to\infty).
\end{equation}
令 $S_n=X_1+\cdots+X_n$，$\mu_n=\E(X_1\one_{\{\abs{X_1}\leqslant n\}})$。则 $S_n/n-\mu_n\toP0$。
\end{theorem}

\begin{remark}[证明思路]
用定理~\ref{thm:weak-laws:truncated-triangular}，取 $X_{n,k}=X_k$（每一行都是同一个序列的前 $n$ 项），$b_n=n$。条件 (i) 就是 $n\P(\abs{X_1}>n)\to0$，这正是假设 \eqref{eq:weak-laws:xP}。条件 (ii) 是 $\frac1n\E(X_1^2;\abs{X_1}\leqslant n)\to0$，要用尾概率公式把它也化成关于 $y\P(\abs{X_1}>y)$ 的条件：它是函数 $2y\P(\abs{X_1}>y)$ 在 $[0,n]$ 上的平均值，而这个函数趋于零。
\end{remark}

\begin{proof}
取 $X_{n,k}=X_k$，$b_n=n$，$\bar X_{n,k}=X_k\one_{\{\abs{X_k}\leqslant n\}}$。

\emph{条件 (i)。}由同分布，$\sum_{k=1}^n\P(\abs{X_k}>n)=n\P(\abs{X_1}>n)\to0$。

\emph{条件 (ii)。}由同分布，要证 $n^{-2}\cdot n\E\bar X_{n,1}^2=\frac1n\E\bar X_{n,1}^2\to0$。由尾概率公式（定理~\ref{thm:product:tail-formula}(2)，$p=2$），
\[
\E\bar X_{n,1}^2=\int_0^\infty2y\P(\abs{\bar X_{n,1}}>y)\dd y\leqslant\int_0^n2y\P(\abs{X_1}>y)\dd y .
\]
这里用了：$y\geqslant n$ 时 $\P(\abs{\bar X_{n,1}}>y)=0$（$\abs{\bar X_{n,1}}\leqslant n$）；$y<n$ 时 $\{\abs{\bar X_{n,1}}>y\}=\{y<\abs{X_1}\leqslant n\}\subseteq\{\abs{X_1}>y\}$。

令 $g(y)=2y\P(\abs{X_1}>y)$，$y\geqslant0$。它是 Borel 函数（$y$ 与单调函数之积），$0\leqslant g(y)\leqslant2y$，并且由 \eqref{eq:weak-laws:xP}，$g(y)\to0$（$y\to\infty$）。所以 $g$ 有界：取 $Y$ 使 $y>Y$ 时 $g(y)\leqslant1$，则 $M=\sup_yg(y)\leqslant\max\{2Y,1\}<\infty$。作变量替换 $y=nx$（第~\ref{lec:integration} 讲作业 E），
\[
\frac1n\int_0^ng(y)\dd y=\int_0^1g(nx)\dd x .
\]
对每个 $x\in(0,1]$，$g(nx)\to0$；$\abs{g(nx)}\leqslant M$，常数在 $[0,1]$ 上可积。由有界收敛定理（推论~\ref{cor:integration:bct}），右边趋于 $0$。

由定理~\ref{thm:weak-laws:truncated-triangular}，$(S_n-n\mu_n)/n\toP0$，即 $S_n/n-\mu_n\toP0$。
\end{proof}

\begin{remark}
(1) 条件 \eqref{eq:weak-laws:xP} 是必要的：若存在常数 $a_n$ 使 $S_n/n-a_n\toP0$，则 \eqref{eq:weak-laws:xP} 成立。这个方向的证明较长，我们不证，可参看 Feller 的书（第二卷）。

(2) \eqref{eq:weak-laws:xP} 只比 $\E\abs{X_1}<\infty$ 弱一点点。由尾概率公式（$p=1-\varepsilon$），若 $x\P(\abs{X_1}>x)\to0$，则对 $x\geqslant x_0$ 有 $\P(\abs{X_1}>x)\leqslant\frac1x$，于是
\[
\E\abs{X_1}^{1-\varepsilon}=\int_0^\infty(1-\varepsilon)x^{-\varepsilon}\P(\abs{X_1}>x)\dd x\leqslant\int_0^{x_0}(1-\varepsilon)x^{-\varepsilon}\dd x+\int_{x_0}^\infty(1-\varepsilon)x^{-1-\varepsilon}\dd x<\infty .
\]
反过来，$\E\abs{X_1}<\infty$ 蕴含 \eqref{eq:weak-laws:xP}（第~\ref{lec:expectation} 讲的例题，第~\pageref{we:expectation:tP} 页）。$\E\abs{X_1}=\infty$ 而 \eqref{eq:weak-laws:xP} 成立的例子见作业 B、C。
\end{remark}

最后得到弱大数定律最常见的形式。

\begin{theorem}[Khinchin 弱大数定律]\label{thm:weak-laws:khinchin}
设 $X_1,X_2,\dots$ 是 i.i.d.，$\E\abs{X_1}<\infty$，$\mu=\E X_1$，$S_n=X_1+\cdots+X_n$。则 $S_n/n\toP\mu$。
\end{theorem}
\begin{proof}
由控制收敛定理用两次：
\begin{itemize}
\item $x\P(\abs{X_1}>x)\leqslant \E(\abs{X_1};\abs{X_1}>x)\to0$（$x\to\infty$），这是第~\ref{lec:expectation} 讲的例题（第~\pageref{we:expectation:tP} 页）；
\item $\mu_n=\E(X_1\one_{\{\abs{X_1}\leqslant n\}})\to \E X_1=\mu$：被积函数逐点趋于 $X_1$，被可积的 $\abs{X_1}$ 控制。
\end{itemize}
由定理~\ref{thm:weak-laws:wlln}，对 $\varepsilon>0$，$\P(\abs{S_n/n-\mu_n}>\varepsilon/2)\to0$。取 $N$ 使 $n\geqslant N$ 时 $\abs{\mu_n-\mu}<\varepsilon/2$，则 $\{\abs{S_n/n-\mu}>\varepsilon\}\subseteq\{\abs{S_n/n-\mu_n}>\varepsilon/2\}$，所以 $\P(\abs{S_n/n-\mu}>\varepsilon)\to0$。
\end{proof}

\begin{example}[Cauchy 分布：弱大数定律不成立]\label{ex:weak-laws:cauchy}
设 $X_1,X_2,\dots$ 独立，都服从 Cauchy 分布，密度 $\frac1{\pi(1+x^2)}$。当 $x\to\infty$ 时，
\[
\P(\abs{X_1}>x)=2\int_x^\infty\frac{\dd t}{\pi(1+t^2)}=\frac2\pi\Bigl(\frac\pi2-\arctan x\Bigr)\sim\frac2\pi\int_x^\infty\frac{\dd t}{t^2}=\frac2{\pi x},
\]
所以 $x\P(\abs{X_1}>x)\to\frac2\pi\neq0$。由上面注记中的必要性，不存在常数 $\mu_n$ 使 $S_n/n-\mu_n\toP0$。以后学了特征函数，我们会证明一个更惊人的事实：$S_n/n$ 与 $X_1$ 同分布，不论 $n$ 多大。平均了 $n$ 次，一点也没有变得更集中。这是因为 Cauchy 分布的尾部太重：$n$ 个样本中最大的那个大约是 $n$ 的量级，它一个就足以左右平均值。
\end{example}

下面的例子说明，期望为无穷时，弱大数定律在适当的形式下仍可能成立。

\begin{workedexample}[St.\ Petersburg 悖论]\label{we:weak-laws:st-petersburg}
设 $X_1,X_2,\dots$ 独立，$\P(X_i=2^j)=2^{-j}$，$j\geqslant1$。（抛一枚公平的硬币直到出现正面，若用了 $j$ 次，你赢 $2^j$ 元。）$\E X_1=\sum_j2^j2^{-j}=\infty$。证明 $S_n/(n\log_2n)\toP1$。由此说明，玩 $n$ 次这个游戏，每次的“公平价格”大约是 $\log_2n$ 元。
\end{workedexample}
\begin{solution}
本例中的对数都以 $2$ 为底。用定理~\ref{thm:weak-laws:truncated-triangular}，$X_{n,k}=X_k$。关键是选 $b_n$：要使 (i) 成立，$b_n$ 不能太小；要使 (ii) 成立，$b_n$ 要尽量小。

\emph{第一步：尾概率。}对正整数 $m$，$\P(X_1\geqslant2^m)=\sum_{j\geqslant m}2^{-j}=2^{-m+1}$。

\emph{第二步：选 $b_n$。}取正整数 $K(n)\to\infty$，并使 $m(n)=\log n+K(n)$ 是整数（$\log n$ 不是整数时，把 $K(n)$ 取成使和为整数的数），令 $b_n=2^{m(n)}=n2^{K(n)}$。由第一步，
\[
n\P(\abs{X_1}>b_n)\leqslant n\P(X_1\geqslant2^{m(n)})=n2^{-m(n)+1}=2^{-K(n)+1}\to0,
\]
这是条件 (i)。

\emph{第三步：条件 (ii)。}$\bar X_{n,k}=X_k\one_{\{X_k\leqslant b_n\}}$，
\[
\E\bar X_{n,k}^2=\sum_{j=1}^{m(n)}2^{2j}2^{-j}=\sum_{j=1}^{m(n)}2^j\leqslant2^{m(n)+1}=2b_n .
\]
所以 $b_n^{-2}\sum_{k\leqslant n}\E\bar X_{n,k}^2\leqslant\frac{2n}{b_n}=2^{1-K(n)}\to0$。

\emph{第四步：$a_n$。}$\E\bar X_{n,k}=\sum_{j=1}^{m(n)}2^j2^{-j}=m(n)$，所以 $a_n=nm(n)=n\log n+nK(n)$。由定理~\ref{thm:weak-laws:truncated-triangular}，
\[
\frac{S_n-nm(n)}{n2^{K(n)}}\toP0 .
\]

\emph{第五步：选 $K(n)$。}我们希望分母比 $n\log n$ 小，同时 $nK(n)$ 相对于 $n\log n$ 可以忽略。取 $K(n)$ 使 $K(n)\to\infty$，$K(n)\leqslant\log\log n$（$n$ 充分大时；仍要求 $m(n)$ 为整数，这只要 $K(n)$ 在 $[\log\log n-1,\log\log n]$ 中适当选取）。则 $2^{K(n)}\leqslant\log n$，于是
\[
\frac{S_n-nm(n)}{n\log n}=\frac{S_n-nm(n)}{n2^{K(n)}}\cdot\frac{2^{K(n)}}{\log n}\toP0
\]
（第二个因子不超过 $1$）；又 $\frac{nm(n)}{n\log n}=1+\frac{K(n)}{\log n}\to1$。所以 $S_n/(n\log n)\toP1$。

\emph{解释。}期望是无穷，但玩 $n$ 次赢的总钱数大约是 $n\log_2n$，每次大约 $\log_2n$ 元。玩 $1024$ 次，每次的公平价格大约 $10$ 元。Nicolas Bernoulli 在 1713 年写道，没有一个“哪怕只有一半理智的人”会不愿意以 $40$ 个金币的价格卖掉玩这个游戏的权利；按上面的计算，若每次付 $40$ 元，大约要玩 $2^{40}\approx10^{12}$ 次才开始回本。
\end{solution}

\topic{本讲小结}

\begin{enumerate}
\item $L^2$ 弱大数定律（定理~\ref{thm:weak-laws:L2}）：两两不相关、期望相同、方差有界，则 $S_n/n\to\mu$ 在 $L^2$ 中与依概率成立。证明只有三行：不相关时方差可加，$\Var(S_n/n)\leqslant C/n$，再用 Chebyshev 不等式。
\item 三角阵列（定理~\ref{thm:weak-laws:triangular-L2}）：只要 $\Var(S_n)=o(b_n^2)$，就有 $(S_n-\E S_n)/b_n\toP0$。用它解决了集券问题（$T_n\approx n\log n$）、随机置换的圈数（$\approx\log n$）、占位问题（空盒比例 $\approx e^{-c}$）。这些例子中，难处在于把所求的量写成（独立的或协方差可算的）示性函数或简单随机变量之和。
\item 截断法（定理~\ref{thm:weak-laws:truncated-triangular}）：方差无穷时，在适当的水平 $b_n$ 处截断，截掉的部分用“概率小”处理，保留的部分用 Chebyshev 不等式处理。由此得到：i.i.d.\ 且 $x\P(\abs{X_1}>x)\to0$ 时 $S_n/n-\mu_n\toP0$（定理~\ref{thm:weak-laws:wlln}）；$\E\abs{X_1}<\infty$ 时 $S_n/n\toP \E X_1$（Khinchin，定理~\ref{thm:weak-laws:khinchin}）。
\item 反例与边界情形：Cauchy 分布没有弱大数定律；St.\ Petersburg 游戏期望无穷，但 $S_n/(n\log_2n)\toP1$。
\item 弱大数定律只说“对每个大的 $n$，$S_n/n$ 偏离很大的概率很小”，不排除一个固定的 $\omega$ 无穷多次偏离。几乎必然收敛的强大数定律要用到下一讲的 Borel--Cantelli 引理。截断法在那里还会再次出现。
\end{enumerate}

% =====================================================================
\homework{期望无穷时的弱大数定律，一个“不公平的公平游戏”，正随机变量的弱大数定律，Glivenko--Cantelli 定理的依概率形式}[期望无穷时的弱大数定律，不公平的公平游戏，Glivenko–Cantelli 定理]
\label{hw:weak-laws}

\exercisepart{基本习题}

\begin{exercises}{A}{课堂内容的补充}
\item\label{ex:weak-laws:A1} 设 $X_1,X_2,\dots$ 两两不相关，$\E X_i=\mu_i$，$\Var(X_i)/i\to0$（$i\to\infty$）。令 $S_n=X_1+\cdots+X_n$，$\nu_n=\E S_n/n$。证明 $S_n/n-\nu_n\to0$ 在 $L^2$ 中与依概率成立。
\begin{solution}
思路：与定理~\ref{thm:weak-laws:L2} 相同，只是方差不再有界，要用“方差增长得比 $i$ 慢”。

由定理~\ref{thm:weak-laws:variance-sum}，$\E(S_n/n-\nu_n)^2=\Var(S_n/n)=\frac1{n^2}\sum_{i=1}^n\Var(X_i)$。给定 $\varepsilon>0$，取 $N$ 使 $i>N$ 时 $\Var(X_i)\leqslant\varepsilon i$。记 $C_N=\sum_{i\leqslant N}\Var(X_i)$。则对 $n>N$，
\[
\frac1{n^2}\sum_{i=1}^n\Var(X_i)\leqslant\frac{C_N}{n^2}+\frac\varepsilon{n^2}\sum_{i=1}^ni\leqslant\frac{C_N}{n^2}+\varepsilon .
\]
所以 $\limsup_n\E(S_n/n-\nu_n)^2\leqslant\varepsilon$，$\varepsilon$ 任意，$L^2$ 收敛成立。依概率收敛由定理~\ref{thm:convergence:Lp-implies-p}。

这个习题说明，均值可以随 $i$ 变化，方差也可以无界，只要 $\Var(X_i)=o(i)$。
\end{solution}

\item\label{ex:weak-laws:A2} $L^2$ 弱大数定律可以推广到某些不独立的序列。设 $\E X_n=0$，并且对 $m\leqslant n$ 有 $\E(X_nX_m)\leqslant r(n-m)$（左边没有绝对值），其中 $r(k)\to0$（$k\to\infty$）。证明 $(X_1+\cdots+X_n)/n\toP0$。
\begin{solution}
思路：直接计算二阶矩。条件 $m=n$ 时给出 $\E X_n^2\leqslant r(0)<\infty$，所以一切二阶矩有限。

由 $\E S_n^2=\sum_{i=1}^n\E X_i^2+2\sum_{1\leqslant m<k\leqslant n}\E(X_kX_m)$ 与假设，
\[
\E S_n^2\leqslant nr(0)+2\sum_{k=2}^n\sum_{m=1}^{k-1}r(k-m)=nr(0)+2\sum_{k=2}^n\sum_{j=1}^{k-1}r(j)\leqslant2\sum_{k=1}^n\sum_{j=0}^{k-1}r^+(j),
\]
其中 $r^+=\max\{r,0\}$，最后一步把每一项换成不小于它的非负数（$r(0)\geqslant \E X_1^2\geqslant0$）。

由 $r^+(j)\to0$，Cesàro 平均也趋于零：给定 $\varepsilon$，取 $J$ 使 $j\geqslant J$ 时 $r^+(j)\leqslant\varepsilon$，则 $\sum_{j=0}^{k-1}r^+(j)\leqslant R_J+\varepsilon k$，$R_J=\sum_{j<J}r^+(j)$。于是
\[
\E\Bigl(\frac{S_n}n\Bigr)^2\leqslant\frac2{n^2}\sum_{k=1}^n(R_J+\varepsilon k)\leqslant\frac{2R_J}n+\varepsilon\cdot\frac{n+1}n,
\]
$\limsup\leqslant\varepsilon$。所以 $S_n/n\to0$ 在 $L^2$ 中，从而依概率。

例子：平稳序列中相关性随时间间隔衰减（比如 $X_n=\xi_n+\xi_{n+1}$，$\xi_i$ 独立同分布、均值为零），大数定律照样成立。
\end{solution}

\item\label{ex:weak-laws:A3} （Monte Carlo 积分）设 $f$ 是 $[0,1]$ 上的可测函数，$U_1,U_2,\dots$ 独立，都服从 $[0,1]$ 上的均匀分布，$I_n=\frac1n(f(U_1)+\cdots+f(U_n))$，$I=\int_0^1f(x)\dd x$。
\begin{questions}
\item 设 $\int_0^1\abs{f(x)}\dd x<\infty$。证明 $I_n\toP I$。
\begin{solution}
$f(U_1),f(U_2),\dots$ 是独立随机变量的函数，独立（推论~\ref{cor:independence:functions}），同分布。由变量替换公式（推论~\ref{cor:expectation:density}），$\E\abs{f(U_1)}=\int_0^1\abs f<\infty$，$\E f(U_1)=I$。由 Khinchin 弱大数定律（定理~\ref{thm:weak-laws:khinchin}），$I_n\toP I$。
\end{solution}
\item 设 $\int_0^1\abs{f(x)}^2\dd x<\infty$。用 Chebyshev 不等式估计 $\P(\abs{I_n-I}>a/\sqrt n)$。
\begin{solution}
记 $\sigma^2=\Var(f(U_1))=\int_0^1f^2-I^2<\infty$。由定理~\ref{thm:weak-laws:variance-sum}，$\Var(I_n)=\sigma^2/n$。由 Chebyshev 不等式，
\[
\P\bigl(\abs{I_n-I}>a/\sqrt n\bigr)\leqslant\frac{\sigma^2/n}{a^2/n}=\frac{\sigma^2}{a^2}.
\]
所以误差是 $1/\sqrt n$ 的量级：要多得到一位有效数字，样本量要乘以 $100$。这个速度与维数无关，这是 Monte Carlo 方法在高维积分中有用的原因（若 $f$ 定义在 $[0,1]^d$ 上，取 $U_i$ 为 $[0,1]^d$ 上的均匀随机向量，证明完全一样）。
\end{solution}
\end{questions}

\item\label{ex:weak-laws:A4} 设 $X\geqslant0$ 取整数值。
\begin{questions}
\item 证明 $\E X=\sum_{n\geqslant1}\P(X\geqslant n)$。
\begin{solution}
这是命题~\ref{prop:expectation:tail-sum}(1)：逐点 $X=\sum_{n\geqslant1}\one_{\{X\geqslant n\}}$，再逐项求期望。
\end{solution}
\item 求 $\E X^2$ 的类似表达式。
\begin{solution}
对非负整数 $k$，$k^2=\sum_{n=1}^k(2n-1)$（前 $k$ 个奇数之和）。所以逐点 $X^2=\sum_{n\geqslant1}(2n-1)\one_{\{X\geqslant n\}}$（$X=\infty$ 时两边都是 $\infty$），由逐项积分（推论~\ref{cor:integration:series}），
\[
\E X^2=\sum_{n\geqslant1}(2n-1)\P(X\geqslant n).
\]
这与定理~\ref{thm:product:tail-formula}(2) 的公式 $\E X^2=\int_0^\infty2y\P(X>y)\dd y$ 是平行的：在 $y\in[n-1,n)$ 上 $\P(X>y)=\P(X\geqslant n)$，$\int_{n-1}^n2y\dd y=2n-1$。
\end{solution}
\end{questions}

\item\label{ex:weak-laws:A5} 推广尾概率公式：设 $h\geqslant0$ 是 Borel 函数，在有界区间上可积，$H(x)=\int_{(-\infty,x]}h(y)\dd y$（设它对一切 $x$ 有限）。证明对任意随机变量 $X$，
\[
\E H(X)=\int_{-\infty}^\infty h(y)\P(X\geqslant y)\dd y .
\]
重要的特例是 $H(x)=e^{\theta x}$（$\theta>0$）。
\begin{solution}
思路：与定理~\ref{thm:product:tail-formula} 的证明相同，把 $H(X)$ 写成一个关于 $y$ 的积分再交换次序。

对每个 $\omega$，$H(X(\omega))=\int_\R h(y)\one_{\{y\leqslant X(\omega)\}}\dd y$。集合 $G=\{(\omega,y):y\leqslant X(\omega)\}$ 属于 $\cF\otimes\cR$（它是可测函数 $(\omega,y)\mapsto X(\omega)-y$ 的非负集），所以 $(\omega,y)\mapsto h(y)\one_G(\omega,y)$ 是非负可测函数。由 Tonelli 定理（定理~\ref{thm:product:tonelli}），
\[
\E H(X)=\int_\Omega\int_\R h(y)\one_G\dd y\dd \P=\int_\R h(y)\int_\Omega\one_{\{X\geqslant y\}}\dd \P\dd y=\int_\R h(y)\P(X\geqslant y)\dd y .
\]
（两边都可以是 $+\infty$。）特例：$h(y)=\theta e^{\theta y}$，$H(x)=e^{\theta x}$，得 $\E e^{\theta X}=\int_\R\theta e^{\theta y}\P(X\geqslant y)\dd y$。所以 $\E e^{\theta X}<\infty$ 与尾概率 $\P(X\geqslant y)$ 指数衰减的速度有关，这在大偏差估计中常用。
\end{solution}

\item\label{ex:weak-laws:A6} 设 $X_1,X_2,\dots$ 是 i.i.d.，$\E\abs{X_1}<\infty$。证明 $\max_{1\leqslant k\leqslant n}\abs{X_k}/n\toP0$。这说明 $S_n$ 中没有哪一项能左右整体，与 Cauchy 分布的情形（\egref{ex:weak-laws:cauchy}）形成对照。
\begin{solution}
对 $\varepsilon>0$，由次可加性与同分布，
\[
\begin{aligned}
\P\Bigl(\max_{k\leqslant n}\abs{X_k}>n\varepsilon\Bigr)&\leqslant\sum_{k=1}^n\P(\abs{X_k}>n\varepsilon)=n\P(\abs{X_1}>n\varepsilon)\\
&=\frac1\varepsilon\cdot(n\varepsilon)\P(\abs{X_1}>n\varepsilon)\to0,
\end{aligned}
\]
最后一步用了第~\ref{lec:expectation} 讲的例题（第~\pageref{we:expectation:tP} 页）：$t\P(\abs{X_1}>t)\to0$。
\end{solution}

\item\label{ex:weak-laws:A7} （样本方差的相合性）设 $X_1,X_2,\dots$ 是 i.i.d.，$\E X_1^2<\infty$，$\mu=\E X_1$，$\sigma^2=\Var(X_1)$。令 $\bar X_n=\frac1n\sum_{i=1}^nX_i$，$s_n^2=\frac1n\sum_{i=1}^n(X_i-\bar X_n)^2$。证明 $s_n^2\toP\sigma^2$。
\begin{solution}
展开：$s_n^2=\frac1n\sum_iX_i^2-2\bar X_n\cdot\frac1n\sum_iX_i+\bar X_n^2=\frac1n\sum_iX_i^2-\bar X_n^2$。$X_i^2$ 是 i.i.d.，$\E X_1^2<\infty$，由 Khinchin 弱大数定律 $\frac1n\sum_iX_i^2\toP \E X_1^2$；同理 $\bar X_n\toP\mu$。由依概率收敛的运算（推论~\ref{cor:convergence:continuous-mapping}），$\bar X_n^2\toP\mu^2$，$s_n^2\toP \E X_1^2-\mu^2=\sigma^2$。

在统计中，这说明样本方差是总体方差的“相合估计”：样本越大，估计越准。（把 $\frac1n$ 换成 $\frac1{n-1}$ 不影响结论，因为 $\frac n{n-1}\to1$。）
\end{solution}
\end{exercises}

\begin{probleme}{B}{期望无穷，样本均值仍收敛到一个常数}
设 $X_1,X_2,\dots$ 是 i.i.d.，
\[
\P\bigl(X_1=(-1)^kk\bigr)=\frac C{k^2\log k},\qquad k=2,3,\dots,
\]
其中 $C$ 是使概率之和为 $1$ 的常数。本题证明：$\E\abs{X_1}=\infty$，但存在有限的常数 $\mu$ 使 $S_n/n\toP\mu$。
\begin{questions}
\item 证明 $C$ 存在（级数 $\sum_{k\geqslant2}\frac1{k^2\log k}$ 收敛），并证明 $\E\abs{X_1}=\infty$。
\begin{solution}
$k\geqslant3$ 时 $\log k\geqslant1$，$\frac1{k^2\log k}\leqslant\frac1{k^2}$，级数收敛，令 $C=(\sum_{k\geqslant2}\frac1{k^2\log k})^{-1}$。由变量替换公式，
\[
\E\abs{X_1}=\sum_{k\geqslant2}k\cdot\frac C{k^2\log k}=C\sum_{k\geqslant2}\frac1{k\log k}=\infty,
\]
最后一步：$\frac1{x\log x}$ 在 $[2,\infty)$ 上递减，$\sum_{k=3}^n\frac1{k\log k}\geqslant\int_3^{n+1}\frac{\dd x}{x\log x}=\log\log(n+1)-\log\log3\to\infty$。
\end{solution}
\item 证明 $x\P(\abs{X_1}>x)\to0$。
\begin{solution}
对 $x\geqslant3$，
\[
\P(\abs{X_1}>x)=\sum_{k>x}\frac C{k^2\log k}\leqslant\frac C{\log x}\sum_{k>x}\frac1{k^2}\leqslant\frac C{\log x}\int_{x-1}^\infty\frac{\dd t}{t^2}=\frac C{(x-1)\log x},
\]
（第二个不等号：$\frac1{k^2}\leqslant\int_{k-1}^k\frac{\dd t}{t^2}$，而最小的 $k>x$ 满足 $k-1\geqslant x-1$）。所以 $x\P(\abs{X_1}>x)\leqslant\frac{Cx}{(x-1)\log x}\to0$。
\end{solution}
\item 令 $\mu_n=\E(X_1\one_{\{\abs{X_1}\leqslant n\}})$。证明 $\mu_n$ 收敛到一个有限的数 $\mu$。
\begin{solution}
$\mu_n=\sum_{k=2}^n(-1)^kk\cdot\frac C{k^2\log k}=C\sum_{k=2}^n\frac{(-1)^k}{k\log k}$。$\frac1{k\log k}$ 单调递减趋于零，由交错级数的 Leibniz 判别法，级数 $\sum_{k\geqslant2}\frac{(-1)^k}{k\log k}$ 收敛。令 $\mu=C\sum_{k\geqslant2}\frac{(-1)^k}{k\log k}$，则 $\mu_n\to\mu$。

注意这个级数只是条件收敛，不是绝对收敛，这正对应着 $\E\abs{X_1}=\infty$：正负项相消，使截断的期望收敛。
\end{solution}
\item 由定理~\ref{thm:weak-laws:wlln} 完成证明。
\begin{solution}
由 B2) 与定理~\ref{thm:weak-laws:wlln}，$S_n/n-\mu_n\toP0$；由 B3)，$\mu_n\to\mu$。与定理~\ref{thm:weak-laws:khinchin} 的证明最后一段相同，$S_n/n\toP\mu$。

这个例子说明，Khinchin 弱大数定律中的 $\E\abs{X_1}<\infty$ 并不是必要的。（但以后会看到，强大数定律 $S_n/n\to\mu$ a.s.\ 一定要求 $\E\abs{X_1}<\infty$。）
\end{solution}
\end{questions}
\end{probleme}

\begin{probleme}{C}{期望无穷，样本均值以可计算的速度趋于无穷}
设 $X_1,X_2,\dots$ 是 i.i.d.，$\P(X_1>x)=\frac e{x\log x}$（$x\geqslant e$），$\P(X_1\geqslant e)=1$。
\begin{questions}
\item 验证这确实给出了一个分布，并证明 $\E X_1=\infty$。
\begin{solution}
$G(x)=\frac e{x\log x}$ 在 $[e,\infty)$ 上连续、严格递减（$x\log x$ 递增），$G(e)=1$，$G(\infty)=0$。令 $F(x)=0$（$x<e$），$F(x)=1-G(x)$（$x\geqslant e$），它是连续的分布函数，由定理~\ref{thm:distributions:characterization} 存在对应的随机变量。由尾概率公式，
\[
\E X_1=\int_0^\infty \P(X_1>x)\dd x\geqslant\int_e^\infty\frac{e\dd x}{x\log x}=e\bigl[\log\log x\bigr]_e^\infty=\infty .
\]
\end{solution}
\item 证明 $x\P(X_1>x)\to0$。
\begin{solution}
$x\P(X_1>x)=\frac e{\log x}\to0$。
\end{solution}
\item 计算 $\mu_n=\E(X_1\one_{\{X_1\leqslant n\}})$（$n\geqslant e$），并证明 $\mu_n\sim e\log\log n$。
\begin{solution}
由尾概率公式，对 $Y=X_1\one_{\{X_1\leqslant n\}}\geqslant0$，$\P(Y>x)=\P(x<X_1\leqslant n)=\P(X_1>x)-\P(X_1>n)$（$x<n$），$=0$（$x\geqslant n$）。所以
\[
\begin{aligned}
\mu_n&=\int_0^n\bigl(\P(X_1>x)-\P(X_1>n)\bigr)\dd x=e+\int_e^n\frac{e\dd x}{x\log x}-n\P(X_1>n)\\
&=e+e\log\log n-\frac e{\log n},
\end{aligned}
\]
其中 $\int_0^e\P(X_1>x)\dd x=e$。于是 $\mu_n/(e\log\log n)\to1$。
\end{solution}
\item 证明 $S_n/n-\mu_n\toP0$，从而 $\frac{S_n}{n\log\log n}\toP e$。
\begin{solution}
由 C2) 与定理~\ref{thm:weak-laws:wlln}，$S_n/n-\mu_n\toP0$。两边除以 $e\log\log n\to\infty$，仍趋于零；由 C3)，$\mu_n/(e\log\log n)\to1$。所以 $\frac{S_n}{ne\log\log n}\toP1$。

与 St.\ Petersburg 悖论比较：那里尾概率是 $\P(X>x)\approx\frac2x$，$S_n$ 的增长是 $n\log n$；这里尾部轻一点，$\approx\frac e{x\log x}$，$S_n$ 的增长是 $n\log\log n$，只比 $n$ 快一点。
\end{solution}
\end{questions}
\end{probleme}

\thinkingpart

\begin{probleme}[问题]{D}{一个“不公平的公平游戏”}
令 $p_k=\frac1{2^kk(k+1)}$，$k=1,2,\dots$，$p_0=1-\sum_{k\geqslant1}p_k$。设 $X_1,X_2,\dots$ 是 i.i.d.，$\P(X_n=-1)=p_0$，$\P(X_n=2^k-1)=p_k$，$k\geqslant1$。令 $S_n=X_1+\cdots+X_n$，$m(n)=\min\{m:\ 2^{-m}m^{-3/2}\leqslant n^{-1}\}$，$b_n=2^{m(n)}$。
\begin{questions}
\item 证明 $\E X_n=0$。（这是一个“公平”的赌博：每次输 $1$ 元，或者赢 $2^k-1$ 元。）
\begin{solution}
$\sum_{k\geqslant1}2^kp_k=\sum_{k\geqslant1}\frac1{k(k+1)}=\sum_k\bigl(\frac1k-\frac1{k+1}\bigr)=1$。$X_n$ 的期望（各项绝对收敛，因为 $\sum_k2^kp_k<\infty$）是
\[
\E X_n=-p_0+\sum_{k\geqslant1}(2^k-1)p_k=-p_0+1-\sum_{k\geqslant1}p_k=-p_0+1-(1-p_0)=0 .
\]
\end{solution}
\item 证明 $m(n)\to\infty$，$m(n)/\log_2n\to1$，并且 $b_n\log_2n/n\to0$。
\begin{solution}
记 $m=m(n)$。$2^{-m}m^{-3/2}$ 关于 $m$ 严格递减趋于零，所以 $m(n)$ 有定义，随 $n$ 单调不减并趋于无穷。由定义，$2^{-m}m^{-3/2}\leqslant\frac1n<2^{-(m-1)}(m-1)^{-3/2}$（$m\geqslant2$），取以 $2$ 为底的对数：
\[
m-1+\tfrac32\log_2(m-1)<\log_2n\leqslant m+\tfrac32\log_2m,
\]
所以 $\log_2n/m\to1$。又 $b_n=2^m<\frac{2n}{(m-1)^{3/2}}$，于是 $\frac{b_n\log_2n}n<\frac{2\log_2n}{(m-1)^{3/2}}\to0$（$\log_2n\sim m$）。
\end{solution}
\item 验证定理~\ref{thm:weak-laws:truncated-triangular} 的条件 (i)：$n\P(\abs{X_1}>b_n)\to0$。
\begin{solution}
$\abs{X_1}>2^m$ 当且仅当 $X_1=2^k-1$，$k\geqslant m+1$。所以
\[
\P(\abs{X_1}>b_n)=\sum_{k>m}p_k\leqslant\frac1{m^2}\sum_{k>m}2^{-k}=\frac{2^{-m}}{m^2}.
\]
于是 $n\P(\abs{X_1}>b_n)\leqslant n2^{-m}m^{-3/2}\cdot m^{-1/2}\leqslant m^{-1/2}\to0$，用了 $m$ 的定义 $n2^{-m}m^{-3/2}\leqslant1$。
\end{solution}
\item 验证条件 (ii)：$n\E\bar X_{n,1}^2/b_n^2\to0$。（提示：$\sum_{k\leqslant m}2^k/k^2\leqslant C2^m/m^2$。）
\begin{solution}
\emph{提示中的估计。}$\sum_{k\leqslant m/2}\frac{2^k}{k^2}\leqslant\sum_{k\leqslant m/2}2^k\leqslant2^{m/2+1}$；$\sum_{m/2<k\leqslant m}\frac{2^k}{k^2}\leqslant\frac4{m^2}\sum_{k\leqslant m}2^k\leqslant\frac{8\cdot2^m}{m^2}$。而 $2^{m/2+1}\leqslant\frac{2^m}{m^2}$ 对充分大的 $m$ 成立。所以 $\sum_{k\leqslant m}\frac{2^k}{k^2}\leqslant\frac{9\cdot2^m}{m^2}$（$m$ 充分大）。

$\bar X_{n,1}$ 取值 $-1$ 或 $2^k-1$（$k\leqslant m$），所以
\[
\E\bar X_{n,1}^2\leqslant1+\sum_{k\leqslant m}2^{2k}p_k=1+\sum_{k\leqslant m}\frac{2^k}{k(k+1)}\leqslant1+\frac{9\cdot2^m}{m^2}.
\]
于是
\[
\frac{n\E\bar X_{n,1}^2}{b_n^2}\leqslant\frac n{4^m}+\frac{9n}{2^mm^2}\leqslant\frac{m^{3/2}}{2^m}+\frac9{m^{1/2}}\to0,
\]
用了 $n\leqslant2^mm^{3/2}$。
\end{solution}
\item 计算 $a_n=n\E\bar X_{n,1}$，并由定理~\ref{thm:weak-laws:truncated-triangular} 证明 $\dfrac{S_n}{n/\log_2n}\toP-1$。
\begin{solution}
由 D1)，$\E\bar X_{n,1}=\E X_1-\sum_{k>m}(2^k-1)p_k=-\sum_{k>m}2^kp_k+\sum_{k>m}p_k$。而 $\sum_{k>m}2^kp_k=\sum_{k>m}\frac1{k(k+1)}=\frac1{m+1}$，$0\leqslant\sum_{k>m}p_k\leqslant2^{-m}$。所以 $a_n=-\frac n{m+1}+\theta_nn2^{-m}$，$0\leqslant\theta_n\leqslant1$。

由定理~\ref{thm:weak-laws:truncated-triangular}，$(S_n-a_n)/b_n\toP0$。于是
\[
\frac{S_n}{n/\log_2n}=\frac{S_n-a_n}{b_n}\cdot\frac{b_n\log_2n}n+\frac{a_n\log_2n}n .
\]
由 D2)，第一项 $\toP0$；第二项 $=-\frac{\log_2n}{m+1}+\theta_n2^{-m}\log_2n\to-1$（$\log_2n\sim m$，$2^{-m}m\to0$）。
\end{solution}
\item 解释结论。
\begin{solution}
每次赌博的期望是 $0$，并且 $\E\abs{X_1}=p_0+\sum_k(2^k-1)p_k<\infty$，可是赌 $n$ 次以后，几乎肯定已经输掉了大约 $n/\log_2n$ 元。原因是：期望中的“补偿”来自极少发生的巨额收益 $2^k-1$，它们的概率小到在 $n$ 次赌博中几乎不会出现（D3)），而每次输掉的 $1$ 元却是实实在在的。

这与 Khinchin 弱大数定律并不矛盾。那里的结论是 $S_n/n\toP0$，而这里 $S_n/n=\frac{S_n}{n/\log_2n}\cdot\frac1{\log_2n}\toP(-1)\cdot0=0$，两者一致。区别在于精度：弱大数定律只说 $S_n=o(n)$，这里说的是 $S_n\approx-n/\log_2n$，后者恰好是 $o(n)$ 的。所以“期望为零”只保证平均每次的盈亏趋于零，并不保证总盈亏不趋于 $-\infty$。
\end{solution}
\end{questions}
\end{probleme}

\begin{probleme}*[问题]{E}{正随机变量的弱大数定律}
设 $X_1,X_2,\dots$ 是 i.i.d.，$\P(0\leqslant X_i<\infty)=1$，$\P(X_i>x)>0$ 对一切 $x$ 成立，分布函数 $F$ 连续。令 $\mu(s)=\E(X_1;X_1\leqslant s)=\int_{[0,s]}x\dd F(x)$，$\nu(s)=\dfrac{\mu(s)}{s(1-F(s))}$。已知：存在常数 $a_n$ 使 $S_n/a_n\toP1$，当且仅当 $\nu(s)\to\infty$（$s\to\infty$）。本题证明充分性。设 $\nu(s)\to\infty$。
\begin{questions}
\item 证明 $s/\mu(s)$ 在 $\mu(s)>0$ 的 $s$ 上连续，并且 $s\to\infty$ 时 $s/\mu(s)\to\infty$。由此，对充分大的 $n$ 可以取 $b_n$ 使 $n\mu(b_n)=b_n$，并且 $b_n\to\infty$。
\begin{solution}
$F$ 连续，所以 $\mu(s)=\E(X_1;X_1\leqslant s)$ 关于 $s$ 连续（$s_k\to s$ 时 $X_1\one_{\{X_1\leqslant s_k\}}\to X_1\one_{\{X_1\leqslant s\}}$，除了在 $\{X_1=s\}$ 上，而它的概率为零；被 $X_1\one_{\{X_1\leqslant s+1\}}$ 控制，由控制收敛定理）。所以 $s/\mu(s)$ 在 $\mu(s)>0$ 处连续。由 $\nu(s)\to\infty$ 与 $1-F(s)>0$，$\mu(s)>0$ 对充分大的 $s$ 成立。又 $\frac{\mu(s)}s=\E\bigl(\frac{X_1}s;X_1\leqslant s\bigr)\to0$（被积函数不超过 $1$，逐点趋于零，有界收敛定理），所以 $s/\mu(s)\to\infty$。由介值定理，对充分大的 $n$ 存在 $b_n$ 使 $b_n/\mu(b_n)=n$；取最大的这样的 $b_n$（集合有界且闭），则 $b_n\to\infty$：否则 $b_n$ 有界，而 $s/\mu(s)$ 在有界闭区间上有界，与 $b_n/\mu(b_n)=n\to\infty$ 矛盾。
\end{solution}
\item 验证定理~\ref{thm:weak-laws:truncated-triangular} 的条件 (i)。
\begin{solution}
$n\P(X_1>b_n)=n(1-F(b_n))=\frac{b_n}{\mu(b_n)}(1-F(b_n))=\frac1{\nu(b_n)}\to0$。
\end{solution}
\item 验证条件 (ii)。（提示：由尾概率公式 $\E(X_1^2;X_1\leqslant b)\leqslant\int_0^b2y(1-F(y))\dd y$，再用 $1-F(y)=\frac{\mu(y)}{y\nu(y)}$。）
\begin{solution}
与定理~\ref{thm:weak-laws:wlln} 的证明相同，$\E\bar X_{n,1}^2\leqslant\int_0^{b}2y(1-F(y))\dd y$，$b=b_n$。给定 $\varepsilon>0$，取 $y_0$ 使 $y\geqslant y_0$ 时 $\nu(y)\geqslant\frac1\varepsilon$。在 $[0,y_0]$ 上 $2y(1-F(y))\leqslant2y$；在 $[y_0,b]$ 上 $2y(1-F(y))=\frac{2\mu(y)}{\nu(y)}\leqslant2\varepsilon\mu(b)$（$\mu$ 单调）。所以
\[
\E\bar X_{n,1}^2\leqslant y_0^2+2\varepsilon b\mu(b).
\]
于是 $\frac{n\E\bar X_{n,1}^2}{b_n^2}\leqslant\frac{ny_0^2}{b_n^2}+2\varepsilon\frac{n\mu(b_n)}{b_n}=\frac{y_0^2}{b_n\mu(b_n)}+2\varepsilon$，用了 $n=b_n/\mu(b_n)$。第一项趋于零（$b_n\to\infty$，$\mu(b_n)$ 不减），$\limsup\leqslant2\varepsilon$。
\end{solution}
\item 完成证明：$S_n/b_n\toP1$。
\begin{solution}
$a_n=n\E\bar X_{n,1}=n\mu(b_n)=b_n$。由定理~\ref{thm:weak-laws:truncated-triangular}，$(S_n-b_n)/b_n\toP0$，即 $S_n/b_n\toP1$。

St.\ Petersburg 游戏（去掉“$F$ 连续”这个技术性的假设）是这种情形：$\mu(s)\approx\log_2s$，$1-F(s)\approx\frac2s$，$\nu(s)\approx\frac12\log_2s\to\infty$；$n\mu(b_n)=b_n$ 给出 $b_n\approx n\log_2n$，与例题（第~\pageref{we:weak-laws:st-petersburg} 页）的结论一致。
\end{solution}
\end{questions}
\end{probleme}

\exercisepart{课外知识：Glivenko--Cantelli 定理的依概率形式}

\begin{probleme}{F}{经验分布函数一致地逼近真实的分布函数}
设 $X_1,X_2,\dots$ 是 i.i.d.，分布函数为 $F$。令 $F_n(x)=\frac1n\sum_{i=1}^n\one_{\{X_i\leqslant x\}}$，$D_n=\sup_x\abs{F_n(x)-F(x)}$（它是随机变量，见命题~\ref{prop:random-variables:empirical}）。本题证明 $D_n\toP0$。以后讲强大数定律时会把它加强为几乎必然收敛。
\begin{questions}
\item 证明对每个 $x$，$F_n(x)\toP F(x)$，并且 $F_n(x-)\toP F(x-)$，其中 $F_n(x-)=\frac1n\sum_i\one_{\{X_i<x\}}$。
\begin{solution}
$\one_{\{X_i\leqslant x\}}$ 是 i.i.d.\ 的 Bernoulli 随机变量，期望 $F(x)$，方差 $\leqslant\frac14$。由 $L^2$ 弱大数定律（定理~\ref{thm:weak-laws:L2}），$F_n(x)\toP F(x)$。同理 $\one_{\{X_i<x\}}$ 的期望是 $\P(X_1<x)=F(x-)$。
\end{solution}
\item 固定正整数 $k$。令 $x_0=-\infty$，$x_k=+\infty$，$x_j=F^{-1}(j/k)$（$1\leqslant j\leqslant k-1$，$F^{-1}$ 是分位数函数），并约定 $F(-\infty)=F_n(-\infty)=0$，$F(+\infty-)=F_n(+\infty-)=1$。证明 $F(x_j-)-F(x_{j-1})\leqslant\frac1k$，$1\leqslant j\leqslant k$。
\begin{solution}
由引理~\ref{lem:distributions:quantile}，$y<x_j$ 时 $F(y)<\frac jk$，所以 $F(x_j-)\leqslant\frac jk$（$j=k$ 时是 $1=\frac kk$）。又 $F(x_{j-1})\geqslant\frac{j-1}k$（第~\ref{lec:distributions} 讲习题~\exref{ex:distributions:A1} 的第 2 小问：$F(F^{-1}(u))\geqslant u$；$j=1$ 时是 $0$）。相减即得。
\end{solution}
\item 记 $\Delta_n=\max_{0\leqslant j\leqslant k}\max\{\abs{F_n(x_j)-F(x_j)},\abs{F_n(x_j-)-F(x_j-)}\}$（端点处按约定理解）。证明 $D_n\leqslant\Delta_n+\frac1k$。
\begin{solution}
任取 $x\in\R$，设 $x_{j-1}\leqslant x<x_j$（若有若干个 $x_j$ 相等，取使区间非空的那个 $j$）。由 $F_n$、$F$ 单调与 F2)，
\[
F_n(x)\leqslant F_n(x_j-)\leqslant F(x_j-)+\Delta_n\leqslant F(x_{j-1})+\tfrac1k+\Delta_n\leqslant F(x)+\tfrac1k+\Delta_n,
\]
\[
F_n(x)\geqslant F_n(x_{j-1})\geqslant F(x_{j-1})-\Delta_n\geqslant F(x_j-)-\tfrac1k-\Delta_n\geqslant F(x)-\tfrac1k-\Delta_n .
\]
所以 $\abs{F_n(x)-F(x)}\leqslant\Delta_n+\frac1k$，对 $x$ 取上确界即得。

这一步把“不可数个 $x$”化成了“有限个点 $x_j$”，用的是 $F_n$、$F$ 的单调性。这与命题~\ref{prop:random-variables:sup-right-continuous} 中把不可数化为可数是同一种想法，只是这里更进一步，化成了有限个。
\end{solution}
\item 证明 $D_n\toP0$。
\begin{solution}
给定 $\varepsilon>0$，取 $k>\frac2\varepsilon$。由 F3)，$\{D_n>\varepsilon\}\subseteq\{\Delta_n>\frac\varepsilon2\}$，而 $\Delta_n$ 是 $2(k+1)$ 个随机变量 $\abs{F_n(x_j)-F(x_j)}$、$\abs{F_n(x_j-)-F(x_j-)}$ 的最大值，由 F1) 它们都依概率趋于零（端点处恒为零）。由次可加性，
\[
\P(D_n>\varepsilon)\leqslant\sum_{j=0}^k\Bigl[\P\bigl(\abs{F_n(x_j)-F(x_j)}>\tfrac\varepsilon2\bigr)+\P\bigl(\abs{F_n(x_j-)-F(x_j-)}>\tfrac\varepsilon2\bigr)\Bigr]\to0 .
\]
\end{solution}
\item 说明这个定理在统计中的意义。
\begin{solution}
$F$ 是未知的总体分布，$F_n$ 是由数据算出来的。定理说，样本量足够大时，经验分布函数在整条实轴上\emph{一致地}接近真实的分布函数。所以由数据得到的任何“看起来像分布”的量（分位数、直方图……）都在逼近总体的相应的量。这就是从样本推断总体的理论依据，所以 Glivenko--Cantelli 定理有时被称为“统计学的基本定理”。
\end{solution}
\end{questions}
\end{probleme}
